\chapter{Conclusion and Future Enhancements}
\label{chap:conclusion}

\section{Summary of Technical Achievements}
The development of \textbf{``Matsuri Nights --- A Japanese Festival Street''} represents a comprehensive exploration of modern real-time computer graphics. Built completely from the ground up using ISO C++20 and Modern OpenGL 3.3 Core Profile, the project fulfills and surpasses all academic criteria mandated by the CSE4102 curriculum at KUET:
\begin{enumerate}
    \item \textbf{Mathematical Rigor:} Exact implementation of hierarchical TRS transformations, normal matrix inversion, Gram-Schmidt camera orthonormalization, and Bishop frame parallel transport.
    \item \textbf{Complex Moving Hierarchies:} Seamless kinematic motion where child nodes (lantern bodies, magic orbs, pedestrian limbs) evaluate physics relative to moving parent coordinate systems.
    \item \textbf{Advanced Optical Illumination:} A forward rendering pipeline managing 16 active dynamic lights with distance attenuation, spotlight cutoff cones, dynamic sky dome transitions, and 16-sample PCF soft shadow mapping.
    \item \textbf{Dual Ray Tracing Extensions:} Integration of both an interactive real-time GPU Whitted ray tracer operating on full-screen quads and a multi-threaded CPU snapshot ray tracer.
    \item \textbf{Automated Verification:} 312 unit test assertions passing with 100\% success and an automated report figure capture harness ensuring complete reproducibility.
\end{enumerate}

\section{Future Enhancements}
While \textit{Matsuri Nights} delivers exceptional visual fidelity and real-time responsiveness ($>60\,\text{FPS}$), several architectural extensions present exciting avenues for future research:
\begin{itemize}
    \item \textbf{Physically-Based Rendering (PBR):} Migrating from Blinn-Phong empirical reflectance to the microfacet Cook-Torrance BRDF with the GGX normal distribution function, metallic-roughness parameterization, and Image-Based Lighting (IBL) via cubemaps.
    \item \textbf{Normal and Parallax Occlusion Mapping:} Utilizing tangent space normal maps and height fields to simulate 3D relief displacement across stone cobblestones and ceramic roof tiles without increasing polygonal geometric complexity.
    \item \textbf{Volumetric Fog and Light Shafts (God Rays):} Ray-marching participating media through the 2048$\times$2048 shadow depth map to render atmospheric haze and volumetric light shafts emanating from lanterns.
    \item \textbf{GPU Skeletal Skinning:} Implementing dual-quaternion or linear-blend vertex skinning on the vertex shader to drive smoothly deforming character meshes rather than rigid hierarchical limb segments.
\end{itemize}

\begin{thebibliography}{99}
\bibitem{learnopengl} Joey de Vries, \textit{LearnOpenGL: Learn Modern OpenGL Graphics Programming in a Step-by-Step Fashion}, Kendall \& Wulf, 2020.
\bibitem{rtr4} Tomas Akenine-M{\"o}ller, Eric Haines, Naty Hoffman, Angelo Pesce, Michal Iwanicki, and S{\'e}bastien Hillaire, \textit{Real-Time Rendering}, 4th Edition, CRC Press, 2018.
\bibitem{shirley} Steve Marschner and Peter Shirley, \textit{Fundamentals of Computer Graphics}, 5th Edition, CRC Press, 2021.
\bibitem{whitted} Turner Whitted, ``An Improved Illumination Model for Shaded Display,'' \textit{Communications of the ACM}, vol. 23, no. 6, pp. 343--349, 1980.
\bibitem{kaykajiya} Timothy L. Kay and James T. Kajiya, ``Ray Tracing Complex Scenes,'' \textit{ACM SIGGRAPH Computer Graphics}, vol. 20, no. 4, pp. 269--278, 1986.
\bibitem{bishop} Richard L. Bishop, ``There is More Than One Way to Frame a Curve,'' \textit{The American Mathematical Monthly}, vol. 82, no. 3, pp. 246--251, 1975.
\bibitem{opengl33} Khronos Group, \textit{The OpenGL Graphics System: A Specification (Version 3.3 Core Profile)}, March 2010.
\end{thebibliography}
